\documentclass[aspectratio=169]{beamer}

\usetheme{Madrid}
\usecolortheme{default}
\usepackage{booktabs}
\usepackage{array}
\usepackage{siunitx}
\usepackage{amsmath}
\usepackage{xcolor}

\definecolor{newinst}{RGB}{180,30,30}
\definecolor{updatedinst}{RGB}{0,90,160}

\newcommand{\newi}[1]{\textcolor{newinst}{#1}}
\newcommand{\updi}[1]{\textcolor{updatedinst}{#1}}

\title{Energetic Particle Detector Review Update}
\author{Dr Ruth Bamford}
\institute{RAL Space}
%\date{\today}
\date{March 28, 2026}

\begin{document}

\frame{\titlepage}

%------------------------------------------------
\begin{frame}{Orbit-Based View of Detector Developments}
\begin{itemize}
\item Instruments are grouped by measurement environment:
\begin{itemize}
    \item heritage LEO mapping / dosimetry systems
    \item MEO radiation-belt instruments
    \item cislunar / lunar-environment concepts
\end{itemize}
\item Highlighting used in the tables:
\begin{itemize}
    \item \textcolor{newinst}{Red}: not in the original baseline list
    \item \textcolor{updatedinst}{Blue}: baseline instrument with updated / changed performance or status
\end{itemize}
\end{itemize}
\end{frame}

%------------------------------------------------
\begin{frame}{LEO Instruments}
\scriptsize
\begin{tabular}{lcccc}
\toprule
Instrument & Type & Proton (MeV) & Mass & Status \\
\midrule
HMRM & Pixel Si & 0.6--500 & 52 g & flown \\
LUCID & Timepix array & $>$10 & 1.2 kg & flown \\
SATRAM & Timepix & ND & 172--380 g & flown \\
\updi{MIRAM} & Timepix3 + Si & 10--400 & $>$150 g & proposed / development \\
MERiT & SSD + APD & 7--100 & 1.13 kg & flown \\
\newi{RADMON (Aalto-1)} & Si + CsI & 10--80 & 360 g & flown \\
RISEPix & 2$\times$ Timepix stack & ND & $>$1 kg & flown \\
VZLUSAT-1 & Timepix & ND & ND & flown \\
\newi{ALTEA} & Si telescopes & ND & ND & flown \\
\newi{DosTel} & Si telescope & ND & ND & flown \\
TEPC & Gas proportional counter & ND & ND & flown \\
\newi{HEPD} & Si + scintillator & 30--200 & 43 kg & flown \\
\bottomrule
\end{tabular}

\vspace{0.3cm}
\footnotesize
LEO remains dominated by miniaturised, mapping, and dosimetry-oriented systems, especially Timepix-class detectors.
\end{frame}

%------------------------------------------------
\begin{frame}{MEO Instruments}
\scriptsize
\begin{tabular}{lcccc}
\toprule
Instrument & Type & Proton (MeV) & Mass & Status \\
\midrule
EMU & Si telescope & 15--100 & 3.5 kg & flown \\
\newi{Merlin} & Si + SURF & $>$40 & 1.7 kg & flown \\
\newi{CEDEX} & Si-based & ND & ND & flown \\
\newi{SREM (GIOVE-B)} & Multi-detector & 23--450 & 2.6 kg & flown \\
NGRM & Si (PSI) & 2--200 & 1.4 kg & flown \\
\newi{ICARE-NG} & Si coincidence & 27--292 & 2.8 kg & flown \\
CEASE3 & Si stacks & 2--150 & 4 kg & flown \\
\newi{CREDANCE} & Si detectors & ND & ND & flown \\
\newi{REPT} & Si stack & 17--200 & 13.4 kg & flown \\
\bottomrule
\end{tabular}

\vspace{0.3cm}
\footnotesize
MEO systems are primarily radiation-belt instruments with thicker silicon stacks and stronger spectral capability.
\end{frame}

%------------------------------------------------
\begin{frame}{Cislunar / Lunar-Orbit Instruments}
\scriptsize
\begin{tabular}{lcccc}
\toprule
Instrument & Type & Energy Range & Platform & Status \\
\midrule
\newi{CRaTER} & Si pairs + tissue-equivalent plastic & ND & LRO lunar orbit & flown \\
\newi{RADOM} & Mini dosimeter spectrometer & ND & Chandrayaan-1 lunar orbit & flown \\
\newi{MATCH} & DSSD + BGO & MeV range & Chang'E-7 & development \\
\newi{PixPAN} & Magnetic spectrometer + Timepix4 & 0.1--10 GeV & ESA LunPAN & concept \\
\newi{NeuPix} & Hybrid pixel detector & 10--100 MeV & ESA LunPAN & concept \\
\newi{IRADCAL} & Scintillator + SiPM & keV--100 MeV & 2U CubeSat concept & concept \\
\bottomrule
\end{tabular}

\vspace{0.3cm}
\footnotesize
Cislunar developments are increasingly hybrid, with explicit focus on lunar albedo particles, neutrons and gamma-rays.
\end{frame}

%------------------------------------------------
\begin{frame}{Comparison Across Orbit Regimes}
\small
\begin{tabular}{lccc}
\toprule
Orbit Regime & Dominant Detector Style & Main Driver & Typical Trend \\
\midrule
LEO & Pixel / compact Si / dosimetry & mapping, dose, variability & low SWaP \\
MEO & Silicon telescope stacks & trapped belts / spectra & higher fidelity \\
Cislunar & Hybrid pixel + scintillator + spectrometer & SEP, GCR, albedo, neutrons & multi-species capability \\
\bottomrule
\end{tabular}

\vspace{0.5cm}
\begin{itemize}
\item LEO: compact deployment and distributed monitoring
\item MEO: spectral discrimination in radiation belts
\item Cislunar: broad environment, secondaries from regolith, neutron sensitivity
\end{itemize}
\end{frame}

%------------------------------------------------
\begin{frame}{Physics Implications}
\small
\begin{itemize}
\item \textbf{LEO:} compact instruments are often sufficient for dose monitoring and spatial mapping
\item \textbf{MEO:} thicker detector stacks are required to resolve trapped particle populations
\item \textbf{Cislunar:} instrument requirements broaden to include:
\begin{itemize}
    \item SEP protons in the 10--100 MeV range
    \item GCR response at higher energies
    \item secondary neutrons and gamma-rays from regolith
\end{itemize}
\item This is why recent cislunar concepts are increasingly hybrid rather than pure silicon telescope or pure pixel systems
\end{itemize}
\end{frame}

%------------------------------------------------
\begin{frame}{Summary}
\begin{itemize}
\item Grouping instruments by orbit clarifies both technology choice and physics driver
\item \textcolor{newinst}{Red entries} identify instruments absent from the original baseline list
\item \textcolor{updatedinst}{Blue entries} identify baseline instruments whose status / performance has changed
\item LEO remains dominated by compact and pixel-based monitors
\item MEO remains the domain of silicon telescope systems
\item Cislunar developments show the clearest shift toward hybrid architectures
\end{itemize}
\end{frame}

\end{document}
